\chapter{Stand der Technik des bestehenden Systems}
\label{ch:IST-Zustand}
% „Aufbau des Grundsystems der Firma xy"

% aufpassen, dass die logische Reihenfolge eingehalten wird
% aufpassen mit den Überschriften – das Spiel mit Über- und Unterüberschrift
% „Aufbau des Grundsystems der Firma xy“
% OK von der Firma einholen
% „das Experiment wurde bei der Firma xy gemacht“
% „habe die Mitarbeiter von denen genommen“
% könnte auch in die Einleitung und danach gar nicht mehr erwähnen
% „wurde durchgeführt bei der Firma x, die sich mit [Kategorien der Informatik] 


\section{Beschreibung des alten Systems}
- 1 Mio Zeilen Code zu verwalten

\section{Funktionalitäten, Architektur, Nutzungskontext}
...

\section{Identifizierte Schwächen / Probleme}
z.B. Usability-Probleme, technische Einschränkungen, fehlende Features

\section{Zustandsanalyse}
\label{ch:IST-Zustand:sec:Zustandsanalyse}
s. \autoref{fig:Anforderungskatalog_ISO_alt}\\
Ableitung funktionaler und nicht-funktionaler Anforderungen\\
- Was muss das neue System mindestens können?\\
- was soll es besser oder zusätzlich leisten?\\
Quellen der Anforderungen:\\
-> Nutzerfeedback, bestehende Dokumentation, eigene Analyse\\

- die Tabellen aus der ISO nutzen / ausfüllen \cite{DIN_EN_ISO_9241-110} \autoref{fig:Anforderungskatalog_ISO_alt}\\
